\documentclass[11pt]{article}
\usepackage[margin=1in]{geometry}
\usepackage{amsmath}
\usepackage{booktabs}
\usepackage{fancyvrb}
\usepackage{xcolor}
\usepackage{titlesec}
\usepackage{enumitem}
\usepackage{hyperref}
\hypersetup{colorlinks=true, linkcolor=blue, urlcolor=blue}

\titleformat{\section}{\large\bfseries\color{blue!50!black}}{\thesection}{1em}{}
\newcommand{\code}[1]{\texttt{\small #1}}

\title{\textbf{SecureGen / S4:\\ Generation-Time Defenses for LLM-Written Smart Contracts}\\[4pt]
\large A plain-language walkthrough of the pipeline}
\author{AIContractCheck project}
\date{\today}

\begin{document}
\maketitle

\begin{abstract}
\noindent This note explains, in accessible terms, what the SecureGen ``S4''
system does and why. The measurement study behind this project found that
vulnerabilities in LLM-generated Solidity are \emph{intrinsic}: you cannot remove
them by rewording the prompt. S4's answer is to put a lightweight security
``verifier'' \emph{inside} the model's generation loop --- a tiny probe that reads
the model's own internal state token-by-token and steers it away from vulnerable
code as it writes. This document walks through the idea, the training pipeline,
the inference pipeline, and the current experimental results.
\end{abstract}

\section{The problem, in one paragraph}
Developers increasingly paste LLM-generated Solidity straight into production,
where it guards real money. Our 2{,}400-contract, 8-LLM study showed that a large
fraction of that code is vulnerable, and --- crucially --- that ``adversarial''
prompts do \emph{not} meaningfully change the vulnerability rate ($p\approx0.997$).
The danger is not a clever attacker's prompt; it is the ordinary developer who
skips the security pipeline. Asking the model nicely (``write it securely'')
does not fix this. So the fix has to live \emph{in the generation process itself}.

\section{The five strategies}
SecureGen is a modular harness where each defense is a pluggable
\emph{Strategy}, and every strategy is scored by the \emph{same} analyzer
(Slither) so comparisons are fair.

\begin{center}
\begin{tabular}{@{}lll@{}}
\toprule
\textbf{Strategy} & \textbf{Idea} & \textbf{Cost} \\
\midrule
S0 baseline & plain generation (control) & 1 call \\
S1 prompt prefix & prepend a security instruction & 1 call \\
S2 / S2b repair loop & generate $\rightarrow$ Slither $\rightarrow$ repair & several calls \\
S3 checkpoint & write a skeleton first, then fill & several calls \\
\textbf{S4 inline probe} & \textbf{steer decoding with a probe} & 1 call, white-box \\
\bottomrule
\end{tabular}
\end{center}

\noindent S2b already gives a clean, \emph{large} improvement (repairing a
contract against Slither feedback). But S2b needs a full compile-and-Slither pass
per fix --- seconds each, and impossible \emph{while the model is still typing}.
\textbf{S4 is the novel contribution}: it replaces the slow analyzer-in-the-loop
with a fast probe that works mid-generation.

\section{The core insight behind S4}
An ``open'' model (one whose weights we run ourselves --- here
\code{Qwen2.5-Coder-7B}) exposes its internal \emph{activations}: the vectors
flowing through each layer as it generates. We ask a simple question: is
vulnerability already visible in those activations \emph{before} the code is even
finished? If yes, a tiny classifier can read that signal cheaply and warn us in
real time.

\begin{center}
\fbox{\parbox{0.9\textwidth}{\centering
\textbf{TEACHER} (Slither, slow, offline) $\;\rightarrow\;$ labels\\
\textbf{STUDENT} (tiny probe on activations) $\;\rightarrow\;$ fast inline signal\\
\textbf{VERIFIER} (Slither again, on the final code) $\;\rightarrow\;$ fair evaluation}}
\end{center}

\noindent The probe is a cheap, differentiable stand-in for a slow static
analyzer --- and it can ride \emph{inside} decoding, where Slither physically
cannot. That is the whole trick.

\section{Pipeline 1 --- Training the probe (one-time)}
\begin{Verbatim}[frame=single,fontsize=\small]
 Qwen corpus contracts          Slither/Semgrep labels
 dataset/qwen/*.sol (237)       (line numbers parsed from dedup_key #L88-L94)
        |                              |
        v                              v
   +-------------------------------------------------+
   |  extract_activations.py                         |
   |  re-feed each contract through Qwen (on the GPU) |
   |  grab the layer-12 residual stream               |
   |  one vector per LINE (the line's final token)    |
   |  label = 1 if any tool flagged that line         |
   +-----------------------+-------------------------+
                           v
             activations.pt   { X:[N,3584]  y:[N]  groups:[N] }
                           v
   +-------------------------------------------------+
   |  train_probe.py -- tiny LINEAR classifier        |
   |  CONTRACT-LEVEL 80/20 split (no line leakage)    |
   +-----------------------+-------------------------+
                           v
              probe_l12.pt   +   AUROC  <-- the paper result
\end{Verbatim}

\section{Pipeline 2 --- Steered decoding (inference)}
\begin{Verbatim}[frame=single,fontsize=\small]
 prompt --> Qwen decodes token by token (with a KV cache for speed)
                |
                v   every ~16 tokens, at a ';' or '}' boundary:
        read layer-12 hidden state --> probe --> risk in [0,1]
                |
        risk > threshold?
           | no                     | yes
           v                        v
        keep going            BACKTRACK to the last safe point,
                              inject "// checks-effects-interactions..."
                              and re-decode the risky span
                |
                v
        final contract --> Slither (the harness verifies it, fairly)
\end{Verbatim}

\noindent Two engineering details make this runnable and correct:
\begin{itemize}[nosep]
  \item \textbf{KV cache}: without it, every new token re-reads the whole
  sequence ($O(n^2)$); with it, cost is $O(n)$. Verified: cache-on produces
  byte-identical output to cache-off.
  \item \textbf{Backtrack sync}: when we rewind the text, we crop the cache to
  the same point, so the model never ``remembers'' code we erased.
\end{itemize}

\section{How we test ``does it work?''}
For each prompt we run the \emph{same} white-box model twice with the \emph{same}
random seed:
\begin{itemize}[nosep]
  \item \textbf{OFF} --- probe disabled: plain Qwen (the control).
  \item \textbf{ON} --- probe active: steering + backtracking.
\end{itemize}
Because the seed matches, the two runs are identical until the probe first
intervenes, so any difference in Slither's verdict is caused by \emph{steering
alone}. We then report the \textbf{paired $\Delta$vulns}: the per-prompt change
$(\text{vulns}_{\text{ON}} - \text{vulns}_{\text{OFF}})$, averaged over prompts
where both versions compiled. \textbf{Negative = steering removed
vulnerabilities.} Pairing (each contract vs.\ itself) cancels the noise from
contracts having different inherent difficulty.

\section{Results so far}
\subsection*{Probe (the signal exists)}
On the RTX~3090, re-feeding 237 compiled Qwen contracts gave 10{,}287 line
samples (13\% flagged). A \emph{linear} probe recovers Slither's judgment with:
\begin{center}
\begin{tabular}{@{}lcc@{}}
\toprule
 & \textbf{AUROC} & \textbf{note} \\
\midrule
Layer 12 (chosen) & 0.887 & lowest variance; mid-network \\
Best over sweep (L28) & 0.887 & signal is flat across depth \\
Range across all layers & 0.86--0.89 & robust: decodable everywhere \\
\bottomrule
\end{tabular}
\end{center}
\noindent \textbf{Conclusion: vulnerability is \emph{linearly} decodable from
Qwen's own layer-12 representations.} This is the green light for steering ---
the signal the decoder needs is really there.

\subsection*{Steered decoding (does steering help?)}
A paired A/B over 15 diverse prompts (one per contract category), threshold
0.5, gives a \textbf{mixed} result:
\begin{center}
\begin{tabular}{@{}lc@{}}
\toprule
\textbf{Metric (OFF $\rightarrow$ ON)} & \textbf{Value} \\
\midrule
Compile rate & \textbf{40\% $\rightarrow$ 60\%} \quad(rescued 3 contracts) \\
Mean vulns (compiled only) & 0.83 $\rightarrow$ 1.89 \\
Density /100 LOC & 1.25 $\rightarrow$ 3.25 \\
High-severity total & 0 $\rightarrow$ 3 \\
Paired $\Delta$vulns ($n=6$ both compiled) & \textbf{+0.33} \\
\bottomrule
\end{tabular}
\end{center}
\noindent \textbf{Honest reading:} the probe signal is strong, and steering
\emph{improves compilability} (the nudge helps the model finish valid code). But
it does \emph{not} reduce vulnerability density: on contracts that compiled both
ways vulns were flat except one regression, and the rescued contracts compile
\emph{with} vulnerabilities. \textbf{Conclusion: the signal works; naive
comment-injection steering does not yet reduce vulns.} This motivates the next
steps below (calibration via self-generation, a stronger steering intervention,
and a threshold sweep) --- and it is a genuine, publishable negative on the
naive approach, not a dead end.

\section{Status and what remains}
\begin{center}
\begin{tabular}{@{}ll@{}}
\toprule
\textbf{Done} & \textbf{Next} \\
\midrule
Probe de-risk (AUROC 0.89) & Self-generation rewrite (clean number) \\
Layer sweep $\rightarrow$ layer 12 & Scale A/B to all 300 prompts \\
KV-cache decoder (verified) & Threshold sweep for best trade-off \\
Paired A/B harness & Bigger open models (needs the A6000) \\
\bottomrule
\end{tabular}
\end{center}

\noindent \textbf{One-line summary:} training time turns Slither's slow verdicts
into a fast probe; inference time lets that probe steer the model away from
vulnerable code as it writes --- and the probe is accurate enough (AUROC
$\approx 0.89$) for that steering to be well-founded.

\end{document}
